\begin{resumen}

%\avisoLocalizacionArchivo

\begin{comment}
\noindent Aquí va el resumen en español. Debería ocupar entre 1000 y 1500 palabras a modo orientativo.

\warning{Escribe el resumen en estilo periodístico, desde lo más general a lo más específico.  Desde lo más importante a lo menos importante.  Esto será especialmente útil si es necesario hacer un resumen más breve en otro contexto (presentación, herramienta, etc).  Si el resumen está en formato periodístico bastará con seleccionar los primeros párrafos.}
\end{comment}

Este Trabajo Fin de Grado desarrolla un sistema de inteligencia artificial capaz de planificar y explicar maniobras orbitales, que combina agentes de \acrfull{RL} con un \acrfull{LLM} sobre un entorno de simulación física contrastado con datos de misiones reales. Responde a dos dificultades: muchas maniobras reales no admiten una solución analítica ---el achatamiento de los planetas y el rozamiento atmosférico alteran las trayectorias, y a veces son la propia herramienta, como en el aerofrenado---, y las aplicaciones que las modelan exigen una formación especializada. El sistema aprende por sí mismo las maniobras sin fórmula cerrada, reserva los métodos clásicos para los problemas de óptimo conocido y ofrece todo ello mediante una conversación en lenguaje natural, de modo que un usuario no especialista pueda plantear un objetivo y recibir una maniobra calculada y explicada.

El primer pilar es un entorno de simulación propio, construido sobre la biblioteca de astrodinámica Poliastro, que modela la atmósfera de nueve cuerpos del Sistema Solar y propaga las órbitas incorporando el achatamiento y el arrastre. Su validación fue exigente: el modelo reproduce los datos de las sondas con un factor de desviación inferior a 2, coincide con la teoría de la regresión nodal con un error menor del 1\,\% en los cuerpos de achatamiento moderado y supera una batería de 54 pruebas automáticas. Para Urano y Neptuno, sin mediciones in-situ, la validación se apoyó en una reconstrucción del perfil de densidad a partir de la sonda Voyager 2, que llegó a revelar y corregir errores ---el más llamativo, una densidad de referencia de Neptuno equivocada en un factor 100.

Sobre esa base se desarrollaron cinco agentes de RL entrenados con el algoritmo \acrfull{PPO}. El primero aprende la transferencia de Hohmann sin conocer la fórmula, quedando a un 0,03\,\% del óptimo; el segundo la generaliza a cualquier transferencia y cualquier cuerpo mediante la adimensionalización del problema, con un exceso inferior al 1\,\% y sin reentrenar, y una extensión tridimensional añade el cambio de plano manteniéndolo por debajo del 2\,\%. Los otros dos abordan maniobras sin solución analítica: el aerofrenado, que apura la atmósfera hasta el límite seguro de presión dinámica y resuelve 50 de 50 escenarios en cada uno de los siete cuerpos con atmósfera, y el mantenimiento orbital (\textit{station-keeping}), que conserva la órbita dentro de una banda de tolerancia en 30 de 30 escenarios.

El tercer pilar es una capa de orquestación basada en un LLM. Siguiendo el patrón de \textit{tool use}, el modelo no realiza ningún cálculo: interpreta la petición, selecciona y ejecuta la herramienta adecuada ---un agente de RL o un solucionador clásico--- y explica los resultados, anclado de forma estricta a las cifras que estas devuelven. Reúne trece herramientas que cubren todas las maniobras del trabajo y se resolvió, por recomendación de la dirección, con un modelo de lenguaje abierto y gratuito. El resultado integra en una única herramienta el rigor físico y la accesibilidad, y confirma la idea de partida: el aprendizaje automático puede resolver aquello para lo que no existe fórmula, siempre que se apoye en una física contrastada.

\end{resumen}

\begin{abstract}

%\noindent  Write down the abstract in English here. It must be between 1000 and 1500 words long.

This Bachelor's Thesis develops an artificial-intelligence system able to plan and explain orbital maneuvers, combining reinforcement-learning (RL) agents with a large language model (LLM) on a physics-simulation environment validated against real mission data. It addresses two difficulties: many real maneuvers admit no analytical solution (planetary oblateness and atmospheric drag alter the trajectories, and are sometimes the very tool of the maneuver, as in aerobraking) and the applications that model them demand specialized training. The system learns on its own the maneuvers that lack a closed-form formula, reserves classical methods for problems whose optimum is known, and delivers all of it through a natural-language conversation, so that a non-expert user can state a goal and receive a computed and explained maneuver.

The first pillar is a custom simulation environment, built on the astrodynamics library Poliastro, which models the atmosphere of nine Solar System bodies and propagates orbits accounting for oblateness and drag. Its validation was demanding: the model reproduces probe data with a deviation factor below 2, agrees with the theory of nodal regression to within less than 1\,\% for bodies of moderate oblateness, and passes a battery of 54 automated tests. For Uranus and Neptune, lacking in-situ measurements, the validation relied on a reconstruction of the density profile from the Voyager 2 probe, which even revealed and corrected errors (most strikingly, a Neptune reference density wrong by a factor of 100).

On that basis, five RL agents were developed, trained with the PPO algorithm. The first learns the Hohmann transfer without knowing the formula, coming within 0.03\,\% of the optimum; the second generalizes it to any transfer and any body through a non-dimensionalization of the problem, with an excess below 1\,\% and without retraining, and a three-dimensional extension adds the plane change, keeping it below 2\,\%. The other two tackle maneuvers with no analytical solution: aerobraking, which pushes the atmosphere to the safe dynamic-pressure limit and solves 50 of 50 scenarios in each of the seven bodies with an atmosphere, and orbital maintenance (station-keeping), which keeps the orbit within a tolerance band in 30 of 30 scenarios.

The third pillar is an orchestration layer based on an LLM. Following the tool-use pattern, the model performs no computation: it interprets the request, selects and runs the appropriate tool (an RL agent or a classical solver) and explains the results, strictly anchored to the figures they return. It gathers thirteen tools covering every maneuver in the work and was implemented, on the advisors' recommendation, with an open and free language model. The result integrates physical rigor and accessibility into a single tool, and confirms the starting premise: machine learning can solve what has no formula, provided it rests on validated physics.
\end{abstract}